\section{Model}

\subsection{Trust-game architecture}

The simulations use a repeated trust task with a focal active-inference agent interacting with multiple social partners. The primary trust task is a scripted-partner surface designed to isolate the focal agent's inference and deployment mechanism; fully reciprocal active-inference partners are a future extension rather than the evidence base for the present claims.

For each partner, the focal agent maintains one partner-local \texttt{pymdp.Agent}. The partner-local POMDP has three hidden-state factors:

\begin{equation}
 s_t = (s_{\mathrm{type}}, s_{\mathrm{stance}}, s_{\mathrm{own}}).
\end{equation}

The type factor has four states: cooperator, reciprocator, exploiter, and random. The stance factor has three states: trusting, neutral, and hostile. Stance is the partner's disposition toward the focal agent, not the focal agent's attitude. It is hidden from the focal agent and evolves through action-dependent transition dynamics. The own-action factor is deterministic bookkeeping that records the focal agent's previous action so that payoff can be represented as a standard observation likelihood.

The model uses two observation modalities. The first is the partner's action, cooperate or defect, generated from partner type and stance. The second is payoff, generated from own action, type, and stance after marginalizing over the partner action. In the environment, realized payoff is deterministically derived from joint actions. In the generative model, payoff is a likelihood modality over payoff levels, which preserves the incentive structure while staying within standard factorized POMDP conventions.

Partner-local controls are factorized. A singleton partner-control placeholder is used inside each partner-local agent. A stance-control factor drives action-dependent stance transitions during planning, and an own-action factor determines the executed investment/action and payoff likelihood. In agent-choice runs, partner choice is handled outside each partner-local POMDP by evaluating candidate decisions for each partner and routing the final executed action to the selected partner.

\begin{figure}[t]
\centering
\fbox{\begin{minipage}{0.88\linewidth}
\centering
Partner-local \texttt{pymdp.Agent} for each partner\\
$s_{\mathrm{type}} \times s_{\mathrm{stance}} \times s_{\mathrm{own}}$\\
$o_{\mathrm{action}}, o_{\mathrm{payoff}}$\\[0.3em]
external $\beta_k$ tracker $\rightarrow \gamma_k = \gamma_{\mathrm{base}}/\mathbb{E}[\beta_k]$\\
tracked-only lesion: beta updates continue but gamma is fixed
\end{minipage}}
\caption{Partner-specific affective precision architecture. Each partner has a local POMDP for social inference. The beta tracker is external to the POMDP and modulates policy precision rather than adding another sensory modality or hidden-state factor.}
\label{fig:model-schematic}
\end{figure}

\subsection{Observation and transition structure}

The partner-action likelihood maps type and stance to cooperation probability. For example, a cooperator in a trusting stance is highly likely to cooperate, while an exploiter in a hostile stance is unlikely to cooperate. The payoff likelihood uses the standard trust-game payoff table, with payoff levels $-1$, $1$, $3$, and $5$ corresponding to the outcomes of the focal and partner actions.

The type transition is near-identity with small stochastic drift, allowing nonstationarity without making type arbitrarily volatile. The stance transition is action dependent. Cooperating tends to build trust gradually, moving neutral partners toward trusting and hostile partners slowly toward neutral. Defecting can destroy trust quickly, moving trusting partners toward hostility. This action-dependent stance structure is what makes planning depth meaningful in the current architecture.

Preferences are placed over payoff observations rather than directly over partner actions. The agent has no intrinsic preference for observing cooperation or defection by the partner, except insofar as those observations are informative or connected to payoff. Policy priors are uniform in the supported runs unless otherwise specified.

\subsection{External beta tracker}

The beta tracker operates after each partner interaction. For partner $k$, the tracker computes a bounded prediction-error proxy from the probability assigned to the observed partner action before the update:

\begin{equation}
\epsilon_k = 1 - P(o^{\mathrm{action}}_k = o^{\mathrm{obs}}_k \mid q(s_{\mathrm{type}},s_{\mathrm{stance}})).
\end{equation}

The tracker converts this unsigned prediction error into signed affective charge:

\begin{equation}
\phi_k = \alpha(\sigma_0^2 - \epsilon_k^2).
\end{equation}

When squared prediction error is below the baseline surprise scale $\sigma_0^2$, the charge is positive and favors lower beta, which corresponds to higher expected policy precision. When squared prediction error is above baseline, the charge is negative and favors higher beta, corresponding to lower expected policy precision.

The tracker represents $q(\beta_k)$ as a categorical distribution over discrete beta levels. After applying a persistence transition matrix, it forms a pseudo-likelihood from the affective charge and effective precision at each beta level:

\begin{align}
\mathrm{effective\_precision}_\ell &= 1 / \beta_\ell, \\
\log L_\ell(\phi_k) &= \phi_k \cdot \mathrm{effective\_precision}_\ell, \\
q(\beta_k)' &\propto L(\phi_k) \odot T_\beta q(\beta_k).
\end{align}

The expected beta then modulates the partner-local policy precision:

\begin{equation}
\gamma_k = \frac{\gamma_{\mathrm{base}}}{\mathbb{E}_{q(\beta_k)}[\beta_k]}.
\end{equation}

Policy selection follows the standard active-inference softmax over expected free energy, with the partner-local precision applied before policy inference:

\begin{equation}
 q(\pi \mid k) = \sigma(-\gamma_k G(\pi)).
\end{equation}

Low beta therefore means high policy precision and more decisive deployment. High beta means low policy precision and a more diffuse policy posterior.

\subsection{Variants and lesions}

The affect variant updates beta and uses $\gamma_k$ during policy selection. The no-affect or decoupled baseline fixes the effective policy precision to $\gamma_{\mathrm{base}}$. The tracked-only lesion keeps beta updates available but decouples beta from policy selection by forcing $\gamma_k=\gamma_{\mathrm{base}}$. This is the key deployment lesion: the agent can track partner-specific model fitness, but that state cannot sharpen or soften action selection.

The H5 variants perturb beta dynamics. Alexithymia-like variants dampen the charge update so beta moves weakly. Borderline-like variants increase gain or volatility so beta moves too strongly. Depression-like variants bias the initial state or dynamics toward lower expected precision. These labels are computational shorthand only. They do not validate clinical diagnoses, and the manuscript treats them as parameter perturbations of precision dynamics.
